\chapter{Cold-Email Templates by Audience}
\label{ch:cold-email-templates}

The principles in Chapter~\ref{ch:cold-email-principles} apply everywhere; the right template still depends on
who you're writing to. This chapter collects ready-to-adapt templates for professors, recruiters, companies
sourced from investor portfolios, and internship applications — each with the specific tactics that source
emphasized.

\section{Emailing Professors for Research Positions}

Two independent sources converge on nearly the same rules: keep it to roughly 150--180 words; reference the
professor's actual recent work (a specific paper, not a generic compliment); offer to contribute rather than
asking for guidance; send from your institutional email; avoid AI-sounding language; stay flexible to small or
even unpaid tasks at first; follow up once after 3--5 days if there's no response; and contact 7--10 professors,
prioritizing one per target research group rather than mass-emailing a whole department. One source adds a
specific timing recommendation: Tuesday--Thursday, 2--5\,PM.

\begin{template}
Dear Professor [Name],\\[4pt]
I hope this email finds you well. My name is [your name], and I am a [year and program] student in [department]
at [university]. I am writing to kindly inquire whether you would consider me as a [internship/master's
student], as I am preparing to apply to [program name].\\[4pt]
I have carefully studied your work on [research topic], particularly your research on [paper name]. I was
especially interested in how your approach addresses [one specific thing you liked about the paper]. I would be
very grateful for the opportunity to [internship/pursue master's] under your guidance.\\[4pt]
Previously, I completed [previous work experience], where I worked on [what you did there]. Through this
experience, I developed a strong foundation in [areas/topics], and I am eager to further expand my skills in a
[topic] research setting. I have attached my CV and [any other document] for your reference.\\[4pt]
If possible, I would be grateful for a short virtual meeting at your convenience to introduce myself and discuss
my interests in more detail.\\[4pt]
Thank you very much for your time and consideration.\\[4pt]
[Your name]\\{}
[University/college name]
\end{template}

\begin{keytakeaway}
The recommended structure, either way: introduction (name, year, department, university, and purpose) $\to$
research interest (why \emph{this} professor specifically — requires reading 2--3 of their recent papers) $\to$
brief experience/achievements $\to$ a call to action proposing a short virtual meeting.
\end{keytakeaway}

\section{Emailing Recruiters: ``Be Interested, Not Interesting''}

The central principle here, drawn from Dale Carnegie's ``talk to someone about themselves and they will listen
for hours'': a weak cold email tries to impress the recruiter with your own achievements; a strong one shows
genuine curiosity about \emph{their} experience and asks a thoughtful, open-ended question about it.

\subsection{Finding the Right Recruiter on LinkedIn}

\begin{enumerate}
  \item Open the company's LinkedIn page and navigate to its People section.
  \item Search relevant keywords: \texttt{Talent Acquisition}, \texttt{HR}, \texttt{Recruiting}, or \texttt{Early
  Talent} for internship recruiters specifically.
  \item Apply location and function filters to narrow the results.
  \item Inspect individual profiles for recruiters who posted the vacancy, advertised similar roles, or recruit
  for your target department.
\end{enumerate}

If you can't identify the exact recruiter for your role, a recruiter for the relevant department or hiring
category is a reasonable fallback. For finding a work email once you've identified the person, two tools are
cited: \textbf{Apollo} and \textbf{Prospeo} (browser extensions that look up a professional email from a
LinkedIn profile) — see Chapter~\ref{ch:sourcing-tools} for the fuller tool directory.

\subsection{Three Rules for Personalization}

\begin{enumerate}
  \item \textbf{Be genuinely curious.} Don't manufacture interest in someone's background just to look
  personalized — recruiters can often tell.
  \item \textbf{Ask an open-ended question.} Avoid anything answerable with a simple yes/no; ask something that
  invites a story or a perspective.
  \item \textbf{Make the question unique to the person.} If it could be sent unchanged to every recruiter at the
  company, it needs more personalization.
\end{enumerate}

\begin{template}
Hi [Recruiter's Name],\\[4pt]
I hope you're doing well. My name is [Your Name], and we connected on LinkedIn recently. I noticed that you've
[specific observation about their career, experience, or a post you shared], which caught my attention because
[genuine reason it interests you].\\[4pt]
I was curious to learn more about [specific topic related to their experience]. Would you be open to sharing how
you approached [specific question related to their journey or expertise]?\\[4pt]
For context, I'm a [year] student studying [field] and am currently exploring opportunities in [role or area].
I've also applied for [position, if applicable] at [Company] and would love to learn more about the team and the
company's culture.\\[4pt]
If you're available, I'd appreciate the opportunity for a brief 15-minute coffee chat. I'm happy to work around
your schedule.\\[4pt]
Thank you for your time.\\[4pt]
[Your Name] \textbar\ [LinkedIn Profile] \textbar\ [Portfolio or GitHub, if useful]
\end{template}

\textbf{Subject line examples:} \emph{``Coffee Chat Request for Summer 2025 Roles''}; \emph{``Learning About
Your Experience at [Company]''}; \emph{``Question About Your Experience in [Field].''} Avoid anything vague
(``Hello''), exaggerated (``Urgent Opportunity''), or falsely implying an existing relationship.

\textbf{Send time:} the cited recommendation is Tuesday morning, 8--10\,AM in the recipient's local time (citing
HubSpot) — Monday is catch-up-heavy, Friday is wind-down-heavy, after-hours gets overlooked. Treat this as a
starting point to test, not a universal rule; finance and consulting recipients in particular may run on
different schedules.

\section{Sourcing Companies from Investor Portfolios, Then Writing to Them}

Rather than searching job boards directly, this approach starts from \textbf{investor portfolios} — VCs have
already done diligence on which companies are worth watching, so their portfolio pages are a pre-filtered
company list. Cited directories: \textbf{Y Combinator} (\url{ycombinator.com/jobs}), \textbf{Sequoia Capital}
(\url{india.sequoiacap.com/india/companies}), \textbf{Rainmatter} (\url{rainmatter.com}), \textbf{Kalaari
Capital} (\url{kalaari.com/portfolio}), \textbf{Better Capital} (\url{bettercapital.vc/portfolio}). Don't just
collect names — for each company, check whether its product and likely hiring needs actually match your skills
before adding it to your list.

\subsection{Six Principles for the Email Itself}

\begin{enumerate}
  \item \textbf{Keep it short} — one sentence on why you're contacting them, one or two on your qualifications,
  one specific request, your contact info. Cut generic openers like ``I hope this email finds you well'' when
  they add nothing.
  \item \textbf{Write a strong, informative subject line} — e.g.\ \emph{``State-Level Hackathon Winner Interested
  in CRED''} or \emph{``3-Year Experienced Developer Interested in Backend Role at Cars24''} beat \emph{``Quick
  Question''} or \emph{``Hello.''}
  \item \textbf{Demonstrate proof of work} — a short Loom video introducing yourself and one interesting project
  can work; otherwise a portfolio site, GitHub repo, or deployed demo. Explain what the project does, what you
  personally built, the key decisions, and the result.
  \item \textbf{Mention real accomplishments} — ``I built and deployed a machine learning application using
  Python and LightGBM, with source code and a live demo available here'' beats ``I'm passionate about ML and
  have excellent Python skills,'' because the first is inspectable.
  \item \textbf{State a timeline only when it's genuine} — if you have an actual competing offer with a real
  deadline, say so; never invent urgency to manufacture pressure.
  \item \textbf{Include a clear call to action} — ``I would like to interview for the [Frontend Developer
  Internship] position for [season/year]. Could you let me know the appropriate next step?'' beats ``I would
  love to connect sometime.''
\end{enumerate}

\subsection{Four Templates}

\begin{template}
\textbf{Template A --- Direct outreach for a full-time role}\\[4pt]
Hi [Name],\\[2pt]
I'm a former Frontend Developer at [Company A] and [Company B], and I'll be graduating in [Month, Year].\\[2pt]
I'm interested in joining [Target Company] because of my interest in frontend engineering and building polished,
user-friendly web experiences.\\[2pt]
Could you point me toward the appropriate next step for the interview process?\\[2pt]
[Your Name] \textbar\ [LinkedIn] \textbar\ [Portfolio]
\end{template}

\begin{template}
\textbf{Template B --- Highlighting skills and company alignment}\\[4pt]
Dear [Name],\\[2pt]
I'm interested in learning more about the [Role] opportunity at [Company].\\[2pt]
I believe my experience with [Skill 1], [Skill 2], and [Skill 3] could be relevant to your team. Recently, I
have been focused on [Relevant Project or Goal], which aligns with [Specific Company Initiative or Need].\\[2pt]
I'm particularly interested in [Company] because of [Specific Research or Reason for Alignment].\\[2pt]
Would you be available for a brief conversation about the role and the team's requirements?\\[2pt]
[Your Name] \textbar\ [LinkedIn] \textbar\ [Portfolio]
\end{template}

\begin{template}
\textbf{Template C --- Demonstrating initiative through a project}\\[4pt]
Dear [Name],\\[2pt]
I noticed that [Company] is hiring for a [Position] role, and I'm very interested in the opportunity.\\[2pt]
Rather than simply sending my resume, I wanted to share [Project, Demo, or Other Relevant Work], which I built
to demonstrate my ability to work on [Relevant Problem or Skill]. You can explore it here: [Project Link].\\[2pt]
Could you let me know the appropriate next step in the hiring process?\\[2pt]
[Your Name] \textbar\ [LinkedIn]
\end{template}

\begin{template}
\textbf{Template D --- Contacting an alumnus or similar background}\\[4pt]
Dear [Name],\\[2pt]
My name is [Your Name], and I'm currently studying [Program] at [University]. I noticed that you also graduated
from [University] and have since built a career in [Industry or Role].\\[2pt]
I'm interested in pursuing a similar path and would appreciate hearing about your transition into your current
role. Would you be open to a brief 15-minute coffee chat?\\[2pt]
[Your Name] \textbar\ [LinkedIn]
\end{template}

\subsection{Improve Results with A/B Testing}

Change one variable at a time, keep everything else constant, and measure what actually moves:

\begin{longtable}{p{0.18\textwidth} p{0.26\textwidth} p{0.26\textwidth} p{0.22\textwidth}}
\toprule
\textbf{Variable} & \textbf{Version A} & \textbf{Version B} & \textbf{What to measure} \\
\midrule
\endhead
Subject line & Quick Question & Python and ML Projects Relevant to [Company] & Open or reply rate \\
Call to action & Could we connect? & Could we discuss the [Role] for 15 minutes? & Meaningful responses \\
Project evidence & Resume only & Resume plus relevant project demo & Positive replies, interviews \\
Personalization & Generic intro & Specific reference to the company's work & Reply quality, conversion \\
\bottomrule
\end{longtable}

Don't draw conclusions from a handful of emails — response rates are noisy, and an open rate, a reply rate, and
an interview rate are three different things. Follow up roughly a week after the first email if there's no
response; after two or three unanswered emails, it's reasonable to try a different, relevant contact at the same
company rather than re-sending to the same person.

\section{Internship Outreach: A Seven-Part Template}

This structure is built specifically for students without a Tier-1-college pedigree or much formal experience
— it's designed to turn circumstances into a positive narrative rather than an apology.

\begin{template}
\textbf{Internship Outreach Email — structure}\\[4pt]
\textbf{1. Greeting} --- address the recipient politely by name.\\[4pt]
\textbf{2. Introduction} --- your name, institution, degree, current academic year.\\[4pt]
\textbf{3. Internship interest} --- why you're contacting this company, your preferred domain, while remaining
open (e.g.\ ``...in any domain the team finds me a fit,'' which signals flexibility while still naming a
preference).\\[4pt]
\textbf{4. Relevant experience} --- one significant internship, job, research experience, or project, with
responsibilities and relevant skills. If you have none, substitute a strong academic project, independent
experiment, or certification — don't leave this section empty or invent experience.\\[4pt]
\textbf{5. Leadership and achievements} --- positions of responsibility, competitions, extracurriculars,
entrepreneurial initiatives, where genuinely applicable.\\[4pt]
\textbf{6. Personal motivation} --- why the opportunity matters to you specifically, and what makes you a
promising candidate. If you address a non-Tier-1 background, frame it around what you've \emph{done} in
response, not as an excuse.\\[4pt]
\textbf{7. Request and closing} --- express interest in discussing how you could contribute; mention your
attached resume; provide LinkedIn and a phone number.
\end{template}

\subsection{Choosing Who to Contact, and Using LinkedIn Alongside Email}

Don't default to HR. HR handles large volumes of generic applications at many organizations; someone working
directly in your target department (a data analyst, a researcher, a team lead) often understands the role's
actual requirements better and may still be able to introduce you to the right person even if they can't hire
you directly. Exception: some organizations (certain consulting firms, for instance) run centralized recruiting
where HR genuinely is the right contact — read the organization before deciding.

\begin{longtable}{p{0.2\textwidth} p{0.36\textwidth} p{0.36\textwidth}}
\toprule
\textbf{Aspect} & \textbf{Cold email} & \textbf{LinkedIn outreach} \\
\midrule
\endhead
Primary use & Formal introduction, internship inquiry & Networking, guidance, introductions \\
Message length & Can go into more detail & Initial requests should stay brief \\
Personalization & Tailor to company, role, recipient & Reference their professional background \\
Best approach & A well-structured email to a relevant contact & Build connections, start genuine conversations
\\
\bottomrule
\end{longtable}

Understanding LinkedIn's connection degrees helps you prioritize: 1st-degree (direct connections), 2nd-degree
(connected to your connections — often the most useful, since you may share a mutual contact who can introduce
you), and 3rd-degree (further out). A practical habit: \textbf{build connections before you urgently need an
internship}. Students in their first or second year interested in a specific type of company two years out can
start connecting with relevant people now — a relationship built today may become useful later even with no
immediate vacancy, as long as you're connecting with people genuinely relevant to your interests rather than
indiscriminately.

\begin{keytakeaway}
Common mistakes to avoid across all internship outreach: sending identical messages to everyone; contacting only
HR by default; writing excessively about yourself; treating a non-Tier-1 background as the main story rather
than what you did about it; being too rigid about role preference; sending vague requests with no specific next
step; waiting until you urgently need an internship to start networking; and forgetting to include a working
resume link, LinkedIn profile, and contact details.
\end{keytakeaway}

